\addchap{Glossary} \label{ch:glossary}

\begin{description}

\item[adjective phrase (AdjP)] A phrase headed by an adjective that typically functions as a modifier within NPs or as a predicative complement in VPs.

\item[adverb phrase (AdvP)] A phrase headed by an adverb. Functions primarily as modifier of verbs, adjectives, other adverbs, clauses, and PPs. Rarely functions as complement. Accepts degree modifiers like \textit{very}, \textit{too}, \textit{so} but not \textit{right}, \textit{straight}, \textit{just} (in Standard English).

\item[affixal negation] Negation expressed through prefixes (like \emph{un-}, \emph{dis-}, \emph{non-}) or suffixes attached to words rather than through separate negative words.

\item[agreement] The matching of verb form to properties of the subject, particularly in present tense with third-person singular subjects.

\item[allophone] A variant of a phoneme that doesn't change meaning in a given language. For example, the aspirated [t\textsuperscript{h}] in \textit{top} and unaspirated [t] in \textit{stop} are allophones of /t/ in English.

\item[assimilation] A process in connected speech where a sound becomes more like a neighbouring sound. For instance, \textit{ten books} may be pronounced as [tem b\textupsilon ks].

\item[attributive modifier] A modifier that appears within an NP~-- typically before the head N~-- and attributes some property to it.

\item[auxiliary verb] A verb that carries tense, aspect, or modality (e.g., \textit{be}, \textit{have}, \textit{do}, \textit{can}, \textit{will}) and has the NICER properties: Negation, Inversion, Contraction, Ellipsis, and Rebuttal.

\item[clause] A grammatical structure headed by a verb phrase and usually containing a subject.

\item[clausal negation] Negation that applies to an entire clause, subdivided into verbal negation (using \emph{not} with verbs) and non-verbal negation (using \emph{not} with other constituents).

\item[code-switching] Alternating between languages or language varieties within a single conversation, typically reflecting changes in context or social dynamics.

\item[coda] The consonant(s) following the (vowel) nucleus in a syllable. In \textit{cat} /k\ae t/, /t/ is the coda.

\item[complement selection] The more-or-less arbitrary requirement that certain words license specific prepositions as complements (e.g., \textit{depend on}, \textit{interested in}, \textit{proud of}). A major challenge for language learners.

\item[complementation] The systematic way individual words in head function control what can complete the phrase (e.g., an NP, a PP, a clause, etc.), with different words selecting different types of complements.

\item[connected speech] The continuous flow of sounds in natural speech, often involving processes that modify individual word pronunciations.

\item[construction] A conventional pairing of form and meaning that exists at various levels of linguistic structure, from specific phrases to general syntactic patterns.

\item[coordination test] Diagnostic based on the principle that constituents tend to coordinate with like constituents: PPs coordinate with PPs, AdvPs with AdvPs.

\item[counter-example] An example that disproves a proposed rule or definition, being a case where the rule fails. Good for testing linguistic definitions.

\item[defective verb] A verb lacking one or more of the typical verb forms (e.g., \textit{must}, which has no past tense).

\item[definiteness] A grammatical property that indicates whether a noun phrase refers to something specific and identifiable to both speaker and hearer, typically marked by a determiner.

\item[dependent] A phrase that depends on the head. For example, \textit{the} is a dependent in \textit{the cake}.

\item[determiner] A dependent in a noun phrase that marks it as definite or indefinite (e.g., \textit{\emph{a} pen}, \textit{\emph{the} pen}, \textit{\emph{much more} air}, \textit{\emph{almost ten} pens}, \textit{\emph{my} pen}).

\item[determinative] A word of the lexical category that typically specifies the reference and number of noun phrases, including articles (\textit{the}, \textit{a}), demonstratives (\textit{this}, \textit{that}), and quantifiers (\textit{some}, \textit{many}), distinguished from adjectives by their function and distribution.

\item[determinative phrase (DP)] A phrase headed by a determinative that mostly functions as the determiner in an NP, such as \textit{most}, \textit{not many}, \textit{almost every}, and \textit{just about all}.

\item[direct object] A constituent selected by a transitive verb, usually able to be identified through passivization tests (becomes the subject in passive constructions).

\item[downward-entailing context] An environment where inferences flow from general to specific (if \emph{not all students} is true, then \emph{not some students} is also true), typically licensing negative polarity items.

\item[fallacy of monosemy] The mistaken belief that words have (or should have) only one \enquote{true} meaning.

\item[flexible-complement view] The view that prepositions can take various types of complements (NPs, clauses, PPs) or no complement at all~-- paralleling how verbs take different complements. Contrasts with traditional view requiring NP complements.

\item[form-meaning pairing] The connection between how something is said (its linguistic form) and what it means in a particular context. Grammaticality often depends on using established form-meaning pairs within a speech community.

\item[functional overlap] The phenomenon whereby words from different lexical categories can perform similar syntactic functions.

\item[gradability] The property of adjectives, adverbs, and some determinatives and prepositions that allows them to be modified by degree words like \textit{very}, \textit{quite}, or \textit{extremely}, and to form comparative and superlative forms.

\item[grammaticality] The degree to which a construction conforms to the accepted form--meaning pairings of a specific language community, reflecting regularities in syntax, morphology, and socio-cognitive norms, and regardless of how unconventional the meaning is.

\item[head] The word in a phrase that determines the phrase's grammatical type. \textit{Cake} is the head of the NP \textit{the chocolate cake}; \textit{tall} is the head of the AdjP \textit{very tall}.

\item[indirect object] A constituent that precedes the direct object and typically represents the recipient or beneficiary of an action, often paraphraseable with \emph{to} or \emph{for} phrases.

\item[intonation] The rise and fall of voice pitch over phrases and sentences, conveying various linguistic and paralinguistic meanings.

\item[Jespersen's cycle] A historical process describing how negation markers evolve and change in languages, named after Otto Jespersen. English evolved from Old English \emph{ne} through stages of reinforcement to modern contracted forms like \emph{don't}.

\item[lexical category] A category of words (noun, verb, adjective, etc.) having a coherent set of properties.

\item[lexical verb] All non-auxiliary verbs. Lexical verbs typically carry the primary semantic content, contrasting with auxiliary verbs that provide grammatical information.

\item[light verb construction] A multi-word predicate where a semantically light verb combines with a noun to express an action (e.g., \emph{take a walk} vs.\ \emph{walk}).

\item[linking] The connection of sounds across word boundaries in connected speech, as in \textit{an apple}.

\item[many-to-many relationship] The principle that a grammatical category can perform multiple functions and a single function can be performed by multiple categories.

\item[modal auxiliary] A type of auxiliary verb expressing meanings like possibility, necessity, or future time (e.g., \textit{can}, \textit{must}, \textit{will}).

\item[modality] The expression of meanings such as possibility, necessity, permission, and obligation.

\item[modification test] Diagnostic for distinguishing PPs from AdvPs: PPs generally accept \textit{right}, \textit{straight}, \textit{just} as modifiers; AdvPs generally accept \textit{very}, \textit{too}, \textit{so}.

\item[necessary condition] One that must be present for membership in a category. For example, having a head is a necessary condition for being a human, and having a head is a necessary condition for being a phrase.

\item[negative polarity item] A word or phrase (like \emph{any}, \emph{ever}, \emph{yet}, \emph{at all}) that typically appears only in negative contexts or other downward-entailing environments.

\item[NICER properties] The distinctive properties of auxiliary verbs: Negation, Inversion, Contraction, Ellipsis, and Rebuttal.

\item[non-affirmative item] Alternative term for negative polarity item, emphasizing that these items appear in various non-positive contexts, not just overtly negative ones.

\item[normativity] Claims about how language ought to be used, rather than descriptions of how it is used. Related to prescriptive grammar but broader in scope.

\item[noun] A word that can head a noun phrase. Nouns can typically be singular or plural, follow determiners like \textit{the} and \textit{a}, and take AdjP modifiers. They can denote almost anything~-- entities, actions, states, concepts~-- which is why denotation alone can't define them.

\item[noun, count / non-count] A count noun allows counting and plural marking (\textit{1 book}, \textit{2 ideas}); a non-count noun doesn't (*\textit{1 furniture}, *\textit{2 groceries}). The distinction involves both a semantic construal (whether the referent is seen as individuated) and a grammatical marking (whether the noun takes plural and numeral dependents).

\item[NP-complement view] The traditional (but problematic) view that defines prepositions as words that must take NP complements. Creates unnecessary complexity by forcing words like \textit{before} into three different categories.

\item[nucleus] The core of a syllable, typically a vowel. In \textit{cat} /k\ae t/, /\ae/ is the nucleus.

\item[onset] The consonant(s) preceding the nucleus in a syllable. In \textit{cat} /k\ae t/, /k/ is the onset.

\item[aspect] A grammatical system that typically denotes how a situation extends over time, including simple, progressive, perfect, and perfect-progressive combinations.

\item[particle] A preposition functioning as part of a phrasal verb that can appear either before or after an NP object (\textit{pick up the pen}/\textit{pick the pen up}) but must follow pronoun objects (\textit{pick it up}, not *\textit{pick up it}).

\item[participle] A non-tensed verb form, either present participle (\textit{-ing} form) or past participle (typically \textit{-ed} or irregular forms), used in various constructions including progressive and perfect aspects.

\item[perfect aspect] An aspectual form using \textit{have} + past participle that relates a situation to a later reference point, often expressing some relevance to that reference point.

\item[phoneme] The smallest unit of sound that can distinguish meaning in a language. Changing /b/ to /p/ turns \textit{bat} into \textit{pat}.

\item[phonotactics] The rules governing possible sound combinations in a language. English allows /sl/ as an onset (e.g., \textit{slide}) but not as a coda.

\item[phrase] A grammatical structure consisting of a head word and any optional dependents. The head determines the phrase's type: an NP is headed by a noun, a VP by a verb. A phrase can be a single word (\textit{Stop!} is a VP).

\item[phrasal category] A syntactic unit, such as noun phrases (NPs), verb phrases (VPs), or adjective phrases (AdjPs), consisting of a head and any dependents.

\item[phrasal verb] A verb + preposition(s) combination that often expresses idiomatic meaning (e.g., \textit{give up}, \textit{look up}, \textit{take off}). The preposition may function as a particle that can move around NP objects but not pronoun objects.

\item[plural] Denoting more than one, or a number that doesn't fit singular (e.g., zero, fractions).

\item[polarity sensitivity] The phenomenon whereby certain words and phrases are restricted to appearing in either positive or negative linguistic contexts.

\item[positive polarity item] A word or phrase (like \emph{some}, \emph{already}, \emph{a fair chance}) that resists appearing in negative contexts and prefers positive or affirmative environments.

\item[pragmatic negation] Negation indicated by context, tone, or implicature rather than through explicit negative markers in the sentence structure.

\item[predicative complement] A phrase that follows verbs like \textit{be}, \textit{seem}, \textit{become}, \textit{remain}, \textit{look}, and \textit{feel} and that describes or identifies the subject (or, with some verbs, the object).

\item[predicative complement test] Diagnostic showing that PPs (but not AdvPs) can function as complement after verbs that license predicative complements: \textit{The meeting is at noon} but not *\textit{The meeting is quickly}.

\item[preposition fronting] Fronting an entire PP in an interrogative or relative, typically in formal contexts: \textit{About what did you disagree?} versus the more common stranded version \textit{What did you disagree about?}

\item[preposition phrase (PP)] A phrase headed by a preposition. Functions as both modifier and complement in various phrase types. Can modify nouns (unlike AdvPs) and frequently serves as predicative complement after \textit{be}, \textit{seem}, \textit{remain}.

\item[preposition stranding] When a preposition appears without its complement because the complement has been fronted, as in questions (\textit{What are you looking \emph{at}?}) or relative clauses (\textit{the person you spoke \emph{to}}).

\item[progressive aspect] An aspectual form using \textit{be} + \textit{-ing} participle that expresses relevance not just to what continues but to its eventual conclusion.

\item[reduction] The weakening or omission of sounds in connected speech, often in unstressed syllables. For example, \textit{going to} may be pronounced as \textit{gonna}.

\item[resultative construction] A syntactic pattern where a verb is followed by a result phrase describing the outcome of the situation (e.g., \emph{hammer the metal flat}).

\item[rhythm] The timing pattern of stressed and unstressed syllables in speech.

\item[scope of negation] The portion of a sentence that falls under the semantic influence of a negative element, determining what is being negated and affecting the interpretation of modal verbs and other elements.

\item[semantic bleaching] The historical process by which words gradually lose their original concrete meaning and become more grammatical or functional, often occurring in grammaticalization.

\item[serial verb construction] A sequence of verbs that together express a single complex event, typically involving motion and purpose or manner.

\item[shibboleth] A linguistic feature that distinguishes members of one social group from another, often used (consciously or unconsciously) as a test of belonging. From Hebrew \textit{\v{s}ibb\=ole\d{t}}.

\item[sibilant] A type of fricative or affricate consonant characterized by a high-pitched hissing sound, produced when air is forced through a narrow channel between the tongue and the roof of the mouth or teeth.

\item[singular] Denoting exactly one.

\item[Standard English] The most widely accepted variety of a language, particularly in formal, educational, and professional contexts. Not inherently superior to other varieties, but backed by social and institutional power.

\item[Standard Englishes] The various regional and national varieties of Standard English (e.g., Standard Canadian English, Standard British English), acknowledging that what counts as standard varies across regions while maintaining mutual intelligibility.

\item[stress] The relative prominence given to certain syllables or words, through a combination of pitch, length, and loudness.

\item[sub-clausal negation] Negation that applies to constituents smaller than a full clause, such as negated noun phrases or prepositional phrases.

\item[subject] The grammatical function typically performed by the first major constituent in a sentence, controlling verb agreement and identified through tests like tag questions and passivization.

\item[sufficient condition] One that, if present, guarantees membership in a category. For example, being a square is a sufficient condition for being a rectangle, and being an NP is a sufficient condition for being a phrase.

\item[suprasegmental] Features of speech that extend over more than one sound segment, including stress, intonation, and rhythm.

\item[syllable] A unit of pronunciation typically consisting of a vowel (nucleus) with optional consonants before (onset) and after (coda).

\item[syntactic function] The relational grammatical role that a word or phrase plays in sentence structure, including head, subject, object, modifier, complement, and determiner functions.

\item[tense] A grammatical system that typically denotes the time of the situation; English has two tenses (past and present).

\item[verb] A word that can head a verb phrase. Most verbs have a past-tense form, an \textit{-s} form, an \textit{-ing} form, and a plain form. They typically denote actions or states, but this is a tendency, not a definition~-- \textit{existence} is a noun and \textit{resemble} is a verb.

\item[verb phrase] A phrase headed by a verb, potentially including modifiers and complements. Can be modified by adverb phrases but not adjective phrases.

\item[voicing] The vibration of the vocal cords during the production of a speech sound. Voiced sounds (like /b/, /d/, /g/) involve this vibration, while voiceless sounds (like /p/, /t/, /k/) do not.

\end{description}
